\section{Mô hình khử nhiễu DTLN}

DTLN (Dual-signal Transformation LSTM Network) là mô hình khử nhiễu giọng nói
được đề xuất bởi Westhausen và Meyer (2020), thiết kế cho xử lý thời gian thực
trên thiết bị tài nguyên hạn chế.

\begin{table}[htbp]
\centering
\caption{Các thông số chính của mô hình DTLN}
\label{tab:dtln-params}
\begin{tabular}{>{\raggedright\arraybackslash}p{0.45\linewidth} >{\centering\arraybackslash}p{0.45\linewidth}}
\toprule
\rowcolor{headerblue} \textcolor{white}{\textbf{Thông số}} & \textcolor{white}{\textbf{Giá trị}} \\
\midrule
\rowcolor{rowgray} Tần số lấy mẫu & 16.000 Hz, mono \\
Kích thước khung (blockLen) & 512 mẫu (32 ms) \\
\rowcolor{rowgray} Bước dịch khung (block\_shift) & 128 mẫu (8 ms), overlap 75\% \\
Số bin tần số sau STFT & 257 (= blockLen / 2 + 1) \\
\rowcolor{rowgray} Số đơn vị LSTM mỗi lớp & 128 units \\
Số lớp LSTM mỗi giai đoạn & 2 lớp \\
\rowcolor{rowgray} Dropout & 0,25 \\
Số bộ lọc Encoder Conv1D & 256 \\
\rowcolor{rowgray} Hàm kích hoạt mask & Sigmoid {[}0, 1{]} \\
Hàm mất mát & Negative SNR \\
\rowcolor{rowgray} Tổng số tham số & $\approx$ 1,8 triệu \\
Kích thước file model & 3,9 MB (\texttt{DTLN\_vivos\_best.h5}) \\
\rowcolor{rowgray} Độ trễ xử lý & $\approx$ 32 ms (real-time capable) \\
\bottomrule
\end{tabular}
\end{table}

\subsection{Kiến trúc hai giai đoạn}

DTLN xử lý tín hiệu qua hai giai đoạn bổ sung lẫn nhau:

\textbf{Giai đoạn 1 -- Miền tần số (STFT Domain):}
Tín hiệu đầu vào được chia thành các khung 512 mẫu, áp dụng STFT để thu được
257 bin tần số. Phổ biên độ $|X(f,t)|$ được chuẩn hóa qua InstantLayerNormalization,
xử lý qua 2 lớp LSTM (128 units) với Dropout 0,25, rồi qua Dense + Sigmoid để tạo
mask $M_1 \in [0,1]$. Mask được nhân với phổ biên độ gốc, kết hợp lại với phổ pha
và biến đổi ngược (iFFT + Overlap-and-Add) để thu được tín hiệu $y_1(t)$.

\begin{table}[htbp]
\centering
\caption{Cấu trúc Separation Kernel -- Giai đoạn 1}
\label{tab:sep-kernel}
\begin{tabular}{>{\raggedright\arraybackslash}p{0.2\linewidth} >{\centering\arraybackslash}p{0.35\linewidth} >{\raggedright\arraybackslash}p{0.35\linewidth}}
\toprule
\rowcolor{headerblue} \textcolor{white}{\textbf{Thành phần}} & \textcolor{white}{\textbf{Thông số}} & \textcolor{white}{\textbf{Vai trò}} \\
\midrule
\rowcolor{rowgray} LSTM lớp 1 & 128 units, return\_seq=True & Đặc trưng tần số cục bộ \\
Dropout & 0,25 & Giảm overfitting \\
\rowcolor{rowgray} LSTM lớp 2 & 128 units, return\_seq=True & Phụ thuộc thời gian dài hạn \\
Dense & 257 đơn vị & Ánh xạ về không gian bin tần số \\
\rowcolor{rowgray} Sigmoid & Đầu ra {[}0, 1{]} & Tạo mask $M_1(f,t)$ \\
\bottomrule
\end{tabular}
\end{table}

\textbf{Giai đoạn 2 -- Miền đặc trưng (Feature Domain):}
Tín hiệu $y_1(t)$ được chiếu sang không gian 256 chiều qua Conv1D point-wise
(kernel=1, stride=1). Sau InstantLayerNormalization và Separation Kernel thứ hai
(cùng cấu trúc LSTM), Dense + Sigmoid tạo mask $M_2$ trong không gian đặc trưng.
Giải mã qua Conv1D point-wise (256$\rightarrow$512) với causal padding, rồi Overlap-and-Add
tạo tín hiệu sạch $\hat{y}(t)$.

\begin{table}[htbp]
\centering
\caption{So sánh hai giai đoạn xử lý DTLN}
\label{tab:dual-signal}
\begin{tabular}{>{\raggedright\arraybackslash}p{0.2\linewidth} >{\centering\arraybackslash}p{0.33\linewidth} >{\centering\arraybackslash}p{0.33\linewidth}}
\toprule
\rowcolor{headerblue} \textcolor{white}{\textbf{Tiêu chí}} & \textcolor{white}{\textbf{STFT (GĐ 1)}} & \textcolor{white}{\textbf{Conv1D Encoder (GĐ 2)}} \\
\midrule
\rowcolor{rowgray} Loại biến đổi & Cố định (fixed basis) & Học từ dữ liệu \\
Chiều đầu ra & 257 bin tần số & 256 đặc trưng ẩn \\
\rowcolor{rowgray} Ý nghĩa vật lý & Rõ ràng, diễn giải được & Trừu tượng \\
Vai trò & Loại bỏ nhiễu tổng thể & Tinh chỉnh, phục hồi chi tiết \\
\bottomrule
\end{tabular}
\end{table}

\subsection{Hàm mất mát Negative SNR}

Mô hình được tối ưu hóa bằng hàm mất mát Negative SNR:
\begin{equation}
\text{Loss} = -10 \cdot \log_{10}\left(\frac{E[s^2]}{E[(s - \hat{s})^2 + \epsilon]}\right)
\end{equation}
trong đó $s$ là tín hiệu sạch, $\hat{s}$ là tín hiệu ước lượng, $\epsilon = 10^{-7}$.
Hàm mất mát này tối đa hóa tỷ số tín hiệu trên nhiễu (SNR) của đầu ra,
trực tiếp tối ưu chất lượng âm thanh.
